\documentclass[12pt]{article}
\usepackage[utf8]{inputenc}
\usepackage[T1]{fontenc}
\usepackage{hyperref}
\title{Vaiphei Language Resources for K-12 Curriculum Development}
\author{Compiled by Claude Sonnet 5 (claude-sonnet-5) for Andr\'as Kornai}
\date{}
\begin{document}
\maketitle

\section{Introduction}
Vaiphei (ISO 639-3 vap) is a Kuki-Chin (Sino-Tibetan) language spoken mainly in Churachandpur district in southern Manipur, India, with smaller communities in neighbouring states. The 2011 census counted about 43,000 speakers, according to the search summary. It is written in the Latin script.

\section{Primary Educational Resources}
No Vaiphei primers, textbooks or school programmes were found in the searches.

\section{Linguistic and Reference Materials}
\begin{itemize}
\item Wikipedia, Vaiphei language: \url{https://en.wikipedia.org/wiki/Vaiphei_language}
\item Omniglot, Vaiphei: \url{https://omniglot.com/writing/vaiphei.htm}
\end{itemize}
No Vaiphei dictionary was found in the searches.

\section{Religious and Community Resources}
\begin{itemize}
\item Trinitarian Bible Society, Vaiphei Bible translation news (a revision approved for publication; proofreading delayed by unrest in Manipur): \url{https://www.tbsbibles.org/news/651403/Vaiphei-Bible-Translation-.htm}
\item YouVersion, Vaiphei Bible (version 1062): \url{https://www.bible.com/versions/1062}
\item BibliaTodo, Vaiphei Bible (vaipbsi): \url{https://www.bibliatodo.com/sq/bibla/version/vaipbsi}
\item Joshua Project, Vaiphei in India: \url{https://joshuaproject.net/index.php/people_groups/18291/IN}
\end{itemize}
The Bible history reported in the search results: Romans in 1946, the New Testament in 1957, the New Testament with Psalms in 1961, and a complete Bible in 2006.

\section{Digital and Audio Resources}
Online Bible text is available through the YouVersion and BibliaTodo sites listed above. No other audio or video resources were found.

\section{Curriculum Development Context}
The language is spoken by a sizeable community and has a long-established written Bible, which gives a base for orthography and reading texts. No school materials were found. Any curriculum work would need local educators and church and community bodies, whose details were not found in the searches.
\end{document}